% ============================================================================
% Chapter 8: Discussion & Limitations
% ============================================================================
\chapter{Discussion and Limitations}
\label{ch:discussion}

This chapter provides an honest assessment of ZeroCausal's strengths, weaknesses, and the conditions under which it succeeds or fails.

\section{Calibration Sensitivity to Training Contamination}

Our contamination sweep (Section~\ref{sec:contamination}) reveals a critical boundary of our framework: while the PCMCI causal-graph backbone remains stable under sparse training poisoning, the empirical-CDF calibration step of H-CAS is vulnerable. Overall system performance degrades severely from 0.972 to 0.681 under just 1\% training contamination, dropping to 0.658 at 5\% contamination. In contrast, Isolation Forest (0.829) and PyTorch Autoencoder (0.793) degrade much more gracefully due to baseline outlier insulation.

When attack events leak into the conformal calibration queue, the eCDF baseline is inflated, masking subsequent test attacks. Hardening the calibration step---using trimmed eCDFs, self-supervised anomaly filtering, or the Robust Calibration Filter (RCF)---is a critical direction that has been partially addressed in the CausalMLHybrid architecture through the RobustParCorr and SelfHealingCalibration components.

\section{Scalability Limitations on Enterprise Features}

PCMCI causal discovery on the BETH Linux host telemetry dataset ($d = 882$ features, $T = 360$ training windows) requires 7,842 seconds ($\sim$2.18 hours) on a standard CPU. While this discovery phase is a one-time cost, it scales quadratically with feature dimensions ($O(d^2 \tau_{\max} T)$). In high-density enterprise environments where unique system edges can reach thousands of features, this scaling delay represents a major limitation.

Employing a variance-thresholding or mutual-information-based feature selection pre-pass is a mandatory requirement for deployment rather than an optional enhancement, ensuring only active system channels are monitored.

\section{Causal Sufficiency Assumption}

ZeroCausal assumes causal sufficiency---that no unobserved confounders jointly drive the observed provenance features. This is a standard working assumption in the PCMCI framework~\cite{pcmci2019} and in Pearl's structural causal hierarchy. In enterprise networks, unobserved external variables (e.g., domain controller group-policy pushes, NTP synchronization events, or user shift changes) could induce statistical correlations that PCMCI falsely learns as direct causal edges, resulting in localized false alarms.

However, ZeroCausal is \textit{robust to moderate violations} of this assumption for two reasons:
\begin{enumerate}
\item The conformal prediction layer calibrates alarm thresholds empirically on the observed H-CAS score distribution, automatically absorbing any inflation caused by spurious causal edges.
\item The \texttt{AdaptiveWindowDetector} periodically refits the SCM, allowing stale or spurious edges to be discarded when the data distribution shifts.
\end{enumerate}

The practical consequence of confounded edges is elevated (but bounded) false alarms, not missed attacks---a favorable failure mode for a defense system.

\section{Linear Approximation of Count-Based Mechanisms}

The Linear SCM outperforms the Random Forest SCM on our datasets (Table~\ref{tab:performance}), which might appear to undermine the ``causal mechanism'' framing. We clarify: the \textit{causal graph} (which edges exist between provenance features) is the true invariant that ZeroCausal monitors, not the functional form $f_j$.

The linear model serves as a first-order mechanism approximation: on zero-inflated, sparse count data, parent influence is approximately proportional to parent event counts, and the residual distribution captures deviations. The Random Forest overfits rare count bursts in these heavy-tailed distributions, producing noisier residual rankings.

\section{Mimicry Attacks}

If an attacker is aware of the enterprise SCM, they could execute an APT that exactly mimics normal causal dependencies. However, successful mimicry requires the attacker to simultaneously satisfy \textit{all} $d$ causal equations across $\tau_{\max}$ time lags---a $d \times \tau_{\max}$-dimensional constraint. On OpTC ($d = 130$, $\tau_{\max} = 2$), this means matching 260 regression relationships while executing a multi-stage attack.

This is a \textit{deliberate design trade-off}: ZeroCausal sacrifices defense against fully causal-aware adversaries in exchange for zero-label, online operation. In practice, real APT campaigns inject novel process--file edges that have no causal parents in the baseline, triggering structural novelty detection even without mechanism violations.

\section{Operational High-Recall Limitations}

At a 95\% recall target, ZeroCausal exhibits high false positive rates on enterprise data (93.40\% on OpTC). This is expected: catching 95\% of stealthy, multi-stage APT attacks requires pushing the threshold extremely low, flagging benign administrative behaviors. Standard unsupervised baselines perform similarly poorly at this operating point.

ZeroCausal is designed to operate under a user-defined FPR budget (e.g., 5\%). Conformal prediction bounds the empirical alarm rate by this budget ($4.09\% \pm 0.66\%$ across seeds), ensuring a constant, manageable alert volume.

\section{Ethical Considerations}

All real datasets (DARPA OpTC, StreamSpot, BETH) were collected from controlled testbeds or public repositories and do not contain private user information. TC3 and NODLINK are fully simulated. ZeroCausal is designed strictly for defensive monitoring and lacks offensive intrusion capabilities.
